\documentclass{article}
\usepackage[T1]{fontenc}
\usepackage{lmodern}
\usepackage{graphicx}
\usepackage{amsmath,amssymb}
\usepackage[a4paper,margin=2cm]{geometry}
\usepackage[absolute,overlay]{textpos}
\setlength{\TPHorizModule}{1mm}
\setlength{\TPVertModule}{1mm}
\usepackage[indonesian]{babel}
\usepackage{hyperref}
\setlength{\emergencystretch}{2em}
% unit: Halaman 18

\begin{document}

% === Halaman 18 (python/native) ===

Operasi aritmetika pada GF(2m) 

 Misalkan a = (am-1..a1 a0) dan b = (bm-1...b1 b0) GF(2m) 

•

Penjumlahan:  


    a + b = c = (cm-1..c1 c0), yang dalam hal ini ci = (ai + bi) mod 2, c  GF(2m) 

•

Perkalian: a  b = c = (cm-1..c1 c0 ), yang dalam hal ini c adalah sisa pembagian 

polinom a(x)  b(x) dengan irreducible polynomial derajat m,  c  GF(2m).

Rinaldi Munir/II4021 Kriptografi

Note: Sebuah polinom dikatakan tidak dapat direduksi (irreducible polynomial) jika ia 

tidak dapat dinyatakan sebagai perkalian dari dua buah polinom lain  (kecuali perkalian 

dengan 1 dan dirinya sendiri).

Polinom x2 + 1 dan x4 + x + 1 adalah irreducible di dalam GF(2m),

 tetapi polinom  x5 + x2 + x + 1 reducible karena  x5 + x2 + x + 1 = (x5 + x2 + 1) (x2 + 1)

18

\end{document}
